\section{Pairwise character geometry: the Gram identity}
  \label{sec:gram}

  \begin{lemma}[Gram identity]
    \label{lem:gram-identity}
    For any $s, t \in S$, let $G_{st} := \vec{u}_s \cdot \vec{u}_t$ be the Euclidean
    dot product of the corresponding unit axes in $\R^3$.
    Then
    \begin{equation}
      G_{st} \;=\; -\tfrac{1}{2}\,\chi_\rho(st).
      \label{eq:gram-identity}
    \end{equation}
  \end{lemma}

  \begin{proof}
    Using the Pauli product identity
    $({\vec{u}} \cdot \vec{\sigma})({\vec{v}} \cdot \vec{\sigma})
    = (\vec{u} \cdot \vec{v})\,I + i\,(\vec{u} \times \vec{v}) \cdot \vec{\sigma}$,
    we compute
    \begin{align*}
      \rho(s)\,\rho(t)
      &= \bigl(i\,\vec{u}_s \cdot \vec{\sigma}\bigr)
      \bigl(i\,\vec{u}_t \cdot \vec{\sigma}\bigr)
      \;=\; -\bigl(\vec{u}_s \cdot \vec{\sigma}\bigr)
      \bigl(\vec{u}_t \cdot \vec{\sigma}\bigr) \\
      &= -(\vec{u}_s \cdot \vec{u}_t)\,I
      \;-\; i\,(\vec{u}_s \times \vec{u}_t) \cdot \vec{\sigma}.
    \end{align*}
    Taking the trace and using $\Tr(\vec{w} \cdot \vec{\sigma}) = 0$:
    \[
      \chi_\rho(st)
      \;=\; \Tr\bigl(\rho(s)\rho(t)\bigr)
      \;=\; -2\,G_{st}.
      \qedhere
    \]
  \end{proof}

  \begin{remark}
    Lemma~\ref{lem:gram-identity} is the conceptual core of the paper.
    It shows that the \emph{entire axis geometry} --- the pairwise angles between
    neutral generators in $S^2$ --- is readable directly from the character table:
    $G_{st}$ is computable from $\chi_\rho$ alone, without any geometric input.
    In particular, the Gram matrix $(G_{st})_{s,t \in S}$ is an intrinsic spectral
    invariant of the representation~$\rho$.
  \end{remark}
